\documentclass[11pt,a4paper]{jsarticle}
%
\usepackage{amsmath,amssymb}
\usepackage{bm}
\usepackage{graphicx}
\usepackage{here}
\usepackage{ascmac}
\usepackage{listings}
\usepackage{plistings}
\usepackage{url}
%
\setlength{\textwidth}{\fullwidth}
\setlength{\textheight}{40\baselineskip}
\addtolength{\textheight}{\topskip}
\setlength{\voffset}{-0.2in}
\setlength{\topmargin}{0pt}
\setlength{\headheight}{0pt}
\setlength{\headsep}{0pt}
%
\newcommand{\divergence}{\mathrm{div}\,}  %ダイバージェンス
\newcommand{\grad}{\mathrm{grad}\,}  %グラディエント
\newcommand{\rot}{\mathrm{rot}\,}  %ローテーション
%
\lstset{
	%プログラム言語(複数の言語に対応，C,C++も可)
 	language = Python,
 	%背景色と透過度
 	%backgroundcolor={\color[gray]{.90}},
 	%枠外に行った時の自動改行
 	breaklines = true,
 	%自動開業後のインデント量(デフォルトでは20[pt])	
 	breakindent = 10pt,
 	%標準の書体
 	basicstyle = \ttfamily\scriptsize,
 	%basicstyle = {\small}
 	%コメントの書体
 	commentstyle = {\itshape \color[cmyk]{1,0.4,1,0}},
 	%関数名等の色の設定
 	classoffset = 0,
 	%キーワード(int, ifなど)の書体
 	keywordstyle = {\bfseries \color[cmyk]{0,1,0,0}},
 	%""で囲まれたなどの"文字"の書体
 	stringstyle = {\ttfamily \color[rgb]{0,0,1}},
 	%枠 "t"は上に線を記載, "T"は上に二重線を記載
	%他オプション：leftline，topline，bottomline，lines，single，shadowbox
 	frame = TBrl,
 	%frameまでの間隔(行番号とプログラムの間)
 	framesep = 5pt,
 	%行番号の位置
 	numbers = left,
	%行番号の間隔
 	stepnumber = 1,
	%右マージン
 	%xrightmargin=0zw,
 	%左マージン
	%xleftmargin=3zw,
	%行番号の書体
 	numberstyle = \tiny,
	%タブの大きさ
 	tabsize = 4,
 	%キャプションの場所("tb"ならば上下両方に記載)
 	captionpos = t
}
%
\title{TeX書き方のサンプル}
\author{Ryoichi Matsumoto}
\date{\today}
\begin{document}
\maketitle
%
%
\section{はじめに}
Texの書き方のサンプルをまとめる

\section{記号}

\subsection{ギリシャ文字}

\begin{table} [ h ]
\caption{大文字}
	\begin{center}
	\begin{tabular}{| l | l || l | l || l | l || l | l |} \hline
		文字 & コマンド & 文字 & コマンド & 文字 & コマンド & 文字 & コマンド \\ \hline
		$A$ & A & $B$ & B & $\Gamma$ & \verb|\|Gamma & $\Delta$ & \verb|\|Delta \\ \hline
		$E$ & E & $Z$ & Z & $H$ & H & $\Theta$ & \verb|\|Theta \\ \hline
		$I$ & I & $K$ & K & $\Lambda$ & \verb|\|Lambda & $M$ & M \\ \hline
		$N$ & N & $\Xi$ & \verb|\|Xi & $O$ & O & $\Pi$ & \verb|\|Pi \\ \hline
		$P$ & P & $\Sigma$ & \verb|\|Sigma & $T$ & T & $\Upsilon$ & \verb|\|Upsilon \\ \hline
		$\Phi$ & \verb|\|Phi & $X$ & X & $\Psi$ & \verb|\|Psi & $\Omega$ & \verb|\|Omega \\ \hline
	\end{tabular}
	\end{center}
\end{table}

\begin{table} [ h ]
\caption{小文字}
	\begin{center}
	\begin{tabular}{| l | l || l | l || l | l || l | l |} \hline
		文字 & コマンド & 文字 & コマンド & 文字 & コマンド & 文字 & コマンド \\ \hline
		$\alpha$ & \verb|\|alpha & $\beta$ & \verb|\|beta & $\gamma$ & \verb|\|gamma & $\delta$ & \verb|\|delta \\ \hline
		$\epsilon$ & \verb|\|epsilon & $\zeta$ & \verb|\|zeta & $\eta$ & \verb|\|eta & $\theta$ & \verb|\|theta \\ \hline
		$\iota$ & \verb|\|itoa & $\kappa$ & \verb|\|kappa & $\lambda$ & \verb|\|lambda & $\mu$ & \verb|\|mu \\ \hline
		$\nu$ & \verb|\|nu & $\xi$ & \verb|\|xi & $o$ & o & $\pi$ & \verb|\|pi \\ \hline
		$\rho$ & \verb|\|rho & $\sigma$ & \verb|\|sigma & $\tau$ & \verb|\|tau & $\upsilon$ & \verb|\|upsilon \\ \hline
		$\phi$ & \verb|\|phi & $\xi$ & \verb|\|xi & $\psi$ & \verb|\|psi & $\omega$ & \verb|\|omega \\ \hline
	\end{tabular}
	\end{center}
\end{table}

\begin{table} [ h ]
\caption{変体文字}
	\begin{center}
	\begin{tabular}{| l | l || l | l || l | l || l | l |} \hline
		文字 & コマンド & 文字 & コマンド & 文字 & コマンド & 文字 & コマンド \\ \hline
		$\varepsilon$ & \verb|\|varepsilon & $\vartheta$ & \verb|\|vartheta & $\varpi$ & \verb|\|varpi & $\varrho$ & \verb|\|varrho \\ \hline
		$\varsigma$ & \verb|\|varsigma & $\varphi$ & \verb|\|varphi &  &  &  & \\ \hline
	\end{tabular}
	\end{center}
\end{table}


\section{数式}

\subsection{数式}

\begin{lstlisting}
\begin{equation}
	y = ax + b
\end{equation}
\end{lstlisting}

\begin{equation}
	y = ax + b
\end{equation}

\subsection{maxやminの下に文字を入れる}

\begin{lstlisting}
\begin{equation}
	\underset{a}{\rm b}
\end{equation}
\end{lstlisting}

\begin{equation}
	\underset{a}{\rm b}
\end{equation}

\subsection{関数の書体}

\subsubsection{通常の表示}

\begin{lstlisting}
\begin{equation}
	f(x)
\end{equation}
\end{lstlisting}

\begin{equation}
	f(x)
\end{equation}


\subsubsection{斜体のキャンセル}

\begin{lstlisting}
\begin{equation}
	{\rm f}(x)
\end{equation}
\end{lstlisting}

\begin{equation}
	{\rm f}(x)
\end{equation}


\subsubsection{イタリック}

\begin{lstlisting}
\begin{equation}
	{\it f}(x)
\end{equation}
\end{lstlisting}

\begin{equation}
	{\it f}(x)
\end{equation}


\subsection{数学記号}

\begin{table} [ h ]
\caption{等号}
	\begin{center}
	\begin{tabular}{| l | l | l || l | l | l |} \hline
		記号 & コマンド & 意味 & 記号 & コマンド & 意味 \\ \hline
		$=$ & = & イコール & $\neq$ & \verb|\|neq & ノットイコール \\ \hline
		$\sim$ & \verb|\|sim & ニアリーイコール & $\simeq$ & \verb|\|simeq & ニアリーイコール \\ \hline
		$\approx$ & \verb|\|approx & ニアリーイコール & $\fallingdotseq$ & \verb|\|fallingdotseq & ニアリーイコール \\ \hline
		$\risingdotseq$ & \verb|\|risingdotseq & ニアリーイコール & $\equiv$ & \verb|\|equiv & 合同 \\ \hline
	\end{tabular}
	\end{center}
\end{table}

\begin{table} [ h ]
\caption{不等号}
	\begin{center}
	\begin{tabular}{| l | l | l || l | l | l |} \hline
		記号 & コマンド & 意味 & 記号 & コマンド & 意味 \\ \hline
		$>$ & $>$ & 大なり & $<$ & $<$ & 小なり \\ \hline
		$\geq$ & \verb|\|geq & 大なりイコール & $\geqq $ & \verb|\|geqq & 大なりイコール \\ \hline
		$\leq$ & \verb|\|leq & 小なりイコール & $\leqq $ & \verb|\|leqq & 小なりイコール \\ \hline
		$\gg$ & \verb|\|gg & 十分大きい & $\ll$ & \verb|\|ll & 十分小さい \\ \hline
	\end{tabular}
	\end{center}
\end{table}

\begin{table} [ h ]
\caption{演算子}
	\begin{center}
	\begin{tabular}{| l | l | l || l | l | l |} \hline
		記号 & コマンド & 意味 & 記号 & コマンド & 意味 \\ \hline
		$+$ & $+$ & 和 & $-$ & $-$ & 差 \\ \hline
		$\times$ & \verb|\|times & 積，外積 & $\cdot$ & \verb|\|cdot & 積，内積 \\ \hline
		$\otimes$ & \verb|\|otimes & 直積　& $\div$ & \verb|\|div & 商 \\ \hline
		$\pm$ & \verb|\|pm & プラスマイナス & $\mp$ & \verb|\|mp & マイナスプラス \\ \hline
		$*$ & $*$ & 畳み込み & $\propto$ & \verb|\|propto & 比例 \\ \hline
		$\equiv$ & $\verb|\|equiv$ & 定義 & & &  \\ \hline
	\end{tabular}
	\end{center}
\end{table}

\begin{table} [ h ]
\caption{アクセント}
	\begin{center}
	\begin{tabular}{| l | l | l || l | l | l |} \hline
		記号 & コマンド & 意味 & 記号 & コマンド & 意味 \\ \hline
		$\vec{a}$ & \verb|\|vec\{a\} & ベクトル & $\hat{a}$ & \verb|\|hat\{a\} & 推定値 \\ \hline
		$\bar{a}$ & \verb|\|bar\{a\} & 平均 &  &  &  \\ \hline
	\end{tabular}
	\end{center}
\end{table}

\begin{table} [ h ]
\caption{矢印}
	\begin{center}
	\begin{tabular}{| l | l | l || l | l | l |} \hline
		記号 & コマンド & 意味 & 記号 & コマンド & 意味 \\ \hline
		$\leftarrow$ & \verb|\|leftarrow & 左矢印 & $\Leftarrow$ & \verb|\|Leftarrow & 左矢印 \\ \hline
		$\rightarrow$ & \verb|\|rightarrow & 右矢印 & $\Rightarrow$ & \verb|\|Rightarrow & 右矢印 \\ \hline
		$\uparrow$ & \verb|\|uparrow & 上矢印 & $\Uparrow$ & \verb|\|Uparrow & 上矢印 \\ \hline
		$\downarrow$ & \verb|\|downarrow & 下矢印 & $\Downarrow$ & \verb|\|Downarrow & 下矢印 \\ \hline
	\end{tabular}
	\end{center}
\end{table}


\section{図表}

\subsection{図}

\begin{lstlisting}
\begin{figure} [H]
	\begin{center}
		\includegraphics[clip, height=8cm, bb=-20 0 568 278]{figure/sample.png}
		\caption{sample}
		\label{sample_png}
	\end{center}
\end{figure}
\end{lstlisting}

\begin{figure} [H]
	\begin{center}
		\includegraphics[clip, height=8cm, bb=-20 0 568 278]{figure/sample.png}
		\caption{sample}
		\label{sample_png}
	\end{center}
\end{figure}


\subsection{表}

\begin{lstlisting}
\begin{table} [ h ]
\caption{サンプル}
	\begin{center}
	\begin{tabular}{| l | l | l |} \hline
		a & b & c \\ \hline
		d & e & f \\ \hline
		g & h & i \\ \hline
	\end{tabular}
	\end{center}
\end{table}
\end{lstlisting}

\begin{table} [ h ]
\caption{サンプル}
	\begin{center}
	\begin{tabular}{| l | l | l |} \hline
		a & b & c \\ \hline
		d & e & f \\ \hline
		g & h & i \\ \hline
	\end{tabular}
	\end{center}
\end{table}

\section{レイアウト}
...

\section{引用}
参考文献の一覧は\verb|\|thebibliographyを使用する．本文中で参照するには\verb|\|citeを使用する\cite{bibliography}．

\begin{thebibliography}{9}
	\bibitem{bibliography} http://www.latex-cmd.com/struct/thebibliography.html
\end{thebibliography}



\section{参考文献}

\begin{table} [h]
\caption{参考文献}
	\begin{center}
	\begin{tabular} { p{4cm}  p{8cm} } \hline
		タイトル & URL \\ \hline
		Listings & \url{https://mytexpert.osdn.jp/index.php?Listings} \\
		LaTeX入門/各種パッケージの利用 & \url{https://texwiki.texjp.org/?LaTeX入門%2F各種パッケージの利用} \\
		LaTeXコマンド集 & \url{http://www.latex-cmd.com/} \\
		LaTeX入門 & \url{https://medemanabu.net/latex/} \\
		TeXテンプレート & \url{http://hooktail.org/computer/index.php?TeX%A5%C6%A5%F3%A5%D7%A5%EC%A1%BC%A5%C8} \\
		LaTeXで表のセル内改行はtabularx環境を使うと楽 & \url{https://tgnx8810.wordpress.com/2014/11/29/LaTeXで表のセル内改行はtabularx環境を使うと楽/} \\
		数学記号の表 & \url{https://ja.wikipedia.org/wiki/%E6%95%B0%E5%AD%A6%E8%A8%98%E5%8F%B7%E3%81%AE%E8%A1%A8} \\
		max とか min の下に文字を入れる & \url{http://oh-real.jugem.jp/?eid=219} \\
		\hline
	\end{tabular}
	\end{center}
\end{table}


%
%
\end{document}